\documentclass[a4paper]{article}

% Basic packages
\usepackage[fontset=fandol]{ctex}
\usepackage{amsmath, amssymb}
\usepackage{graphicx}
\usepackage[margin=2.5cm]{geometry}
\usepackage[most]{tcolorbox}
\usepackage{etoolbox}
\usepackage{listings}
\usepackage{hyperref}
\usepackage{booktabs}
\usepackage{subcaption}
\usepackage{float}
\usepackage{tikz}
\IfFileExists{pgfplots.sty}{
    \usepackage{pgfplots}
    \pgfplotsset{compat=1.18}
}{}

% Highlight boxes
\newtcolorbox{knowledgebox}[1]{
    enhanced, colback=blue!5!white, colframe=blue!75!black, colbacktitle=blue!75!black,
    coltitle=white, fonttitle=\bfseries, title=#1, attach boxed title to top left={yshift=-2mm, xshift=2mm},
    boxrule=1pt, sharp corners
}

\newtcolorbox{importantbox}[1]{
    enhanced, colback=yellow!10!white, colframe=yellow!80!black, colbacktitle=yellow!80!black,
    coltitle=black, fonttitle=\bfseries, title=#1, sharp corners
}

\newtcolorbox{warningbox}[1]{
    enhanced, colback=red!5!white, colframe=red!75!black, colbacktitle=red!75!black,
    coltitle=white, fonttitle=\bfseries, title=#1, sharp corners
}

\newtcolorbox{dialoguebox}[1]{
    enhanced, breakable, colback=green!4!white, colframe=green!45!black, colbacktitle=green!45!black,
    coltitle=white, fonttitle=\bfseries, title=#1, sharp corners
}

% Listings
\lstset{
    language=Python,
    basicstyle=\ttfamily\small,
    keywordstyle=\color{blue},
    stringstyle=\color{red!60!black},
    commentstyle=\color{green!60!black},
    breaklines=true,
    frame=single,
    numbers=left,
    numberstyle=\tiny\color{gray},
    captionpos=b,
    extendedchars=false
}

% Front-page metadata
\newcommand{\notetitle}{课程笔记：监督微调（SFT for Agents）}
\newcommand{\noteauthors}{Bowen-mantou（整理）}
\newcommand{\notedate}{2026-10-23}
\newcommand{\videochannel}{B 站「常青藤名校课」· CMU 11-768 AI Agents（英文原声字幕版）}
\newcommand{\videopublishdate}{2026-09-09}
\newcommand{\videoduration}{64:59}
\newcommand{\videourl}{https://www.bilibili.com/video/BV1a6Yx62EH4?p=22}
\newcommand{\videocoverpath}{cover.jpg}

\begin{document}

\begin{titlepage}
\centering
\vspace{1.2cm}
{\huge\bfseries \notetitle\par}
\vspace{0.8cm}
{\large \noteauthors\par}
\vspace{0.3cm}
{\large \notedate\par}
\vspace{1.2cm}

\ifdefempty{\videocoverpath}{
{\small 视频封面因版权未收录，见下方视频链接。\par}
}{
\par
}

\vfill
\begin{tcolorbox}[width=0.9\textwidth, colback=black!2!white, colframe=black!60, sharp corners]
\textbf{视频作者/频道}：\videochannel\par
\textbf{发布日期}：\videopublishdate\par
\textbf{视频时长}：\videoduration\par
\textbf{视频链接}：
\ifdefempty{\videourl}{
未填写
}{
\href{\videourl}{\nolinkurl{\videourl}}
}
\par\textbf{PDF制作Skill}：\href{https://github.com/wdkns/wdkns-skills}{\nolinkurl{https://github.com/wdkns/wdkns-skills}}
\end{tcolorbox}
\end{titlepage}

\tableofcontents
\newpage

%% --- 正文内容开始 --- %%

\section{引言：为什么单独讲「训练」}

这一讲是第一讲关于\textbf{训练（training）}的课：前四讲讨论的都是「不改权重、只在推理时做文章」的 agent 增强（上下文、工具、记忆、技能），而从这一讲开始，我们要把学到的经验写进\textbf{模型权重}，并在接下来的三讲里依次展开 SFT、强化学习基础与深度研究。

\subsection{更新 agent 的三条路}

课堂开场先回顾了一个贯穿全课的三分法——想让 agent 变强，可以更新三种东西：

\begin{center}
\begin{tabular}{p{2.6cm}p{4.6cm}p{4.8cm}}
\toprule
载体 & 优点 & 代价 \\
\midrule
上下文窗口 & 保真度高、改动即时生效 & 贵、噪声大、不决定「什么重要」 \\
外部工件 & 可检查、可编辑、可检索、可移植 & 需要诱导、选择与持续维护 \\
模型权重（本讲） & 推理更快、行为改变更彻底，一次训练后每次调用都带着新行为 & 训练慢、过程不透明、与具体模型绑定、单向不可逆 \\
\bottomrule
\end{tabular}
\end{center}

一句话概括三者的时间尺度：上下文与外部工件可以在「秒」级更新，而改权重每次都要付一次训练成本；权重还不能打开看内部发生了什么，也不能把这次改变原样搬到另一个模型上。因此训练这条路换来的是\textbf{快推理、广覆盖、强行为改变}，代价是\textbf{慢、黑盒、模型专属}。

\subsection{从一条 SWE 轨迹看「训练要学什么」}

课堂用一条软件工程（SWE）轨迹做例子：agent 收到一个 issue、拿到一个沙箱里的若干工具（例如 Bash、文件编辑），它一边产生思考、一边发出命令、读取命令返回的 observation，最终由仓库自带的单元测试判定这次运行是否成功。一条轨迹里有四种角色：system、user、assistant，以及 observation；而真正由模型产生的只有 assistant 消息。

那么一条轨迹究竟能教模型什么？课堂列了四类：第一，\textbf{如何写出可被 harness 执行的工具调用}；第二，\textbf{如何在对话变长时继续下去}（因为训练数据里的轨迹本来就是长的）；第三，\textbf{工具结果应该长什么样}——模型要学会「等一个结果」，而不是自己把结果也写出来；第四，数据里附带的其他东西，比如额外的解说、重复的命令，以及在多种可行解法里被选中的那一种。这最后一点是关键：SFT 会把这些「数据里恰好出现的习惯」一起学走，不管它们好不好。

\subsection{本章小结}
改 agent 有三条路：上下文、外部工件、权重。本讲专攻权重：它换来快推理与彻底的行为改变，代价是训练成本、黑盒性和模型绑定。一条轨迹能教工具调用格式、长对话续写、观测格式，也会连数据里的坏习惯一起教——这决定了后面「选数据、过滤数据」的重要性。

\section{SFT 在训练流水线中的位置：冷启动}

\subsection{预训练、中期训练、SFT、RL}

课堂给出一张完整训练流水线的图：\textbf{预训练}（pre-training）用最多 token 让模型学会下一 token 预测；\textbf{中期训练}（mid-training）是预训练与后训练之间的一切（定义各家不一，甚至有人把 SFT 也算进去）；\textbf{SFT} 与 \textbf{RL} 属于后训练（post-training）。一个关键区别是：预训练、中期训练、SFT 都能找到\textbf{固定语料}来训练，而 RL 通常没有固定语料，它训练的是模型\textbf{自己的 rollout}。

课堂以 Kimi K3 的报告为参照给出了一种常见后训练配方：先做 SFT，\textbf{初始化基线的 agent 能力}；再做 RL，\textbf{训练不同推理强度的领域专家}；最后做 \textbf{on-policy distillation}，把专家合并回一个模型。这里的 SFT 就是所谓「冷启动（cold start）」：它决定 RL 从哪个 checkpoint 出发。课堂强调：在 Kimi K3 中，模型是在\textbf{任何 RL 之前}的 SFT 阶段就学会了调用工具、完成长任务。

\subsection{为什么不能跳过 SFT 直接上 RL}

一个自然的反问是：既然 RL 更强，为什么不从 base model 直接 RL？课堂给了 DeepSeek R1.0 的反例：R1.0 在 base model 上直接做 RL、完全跳过 SFT，结果模型的回答\textbf{中英混杂、难以阅读}；在紧接着的下一代模型里，他们又把 SFT 加了回来——R1 从几千条精心挑选的 SFT 样例出发。课堂的概括是：\textbf{SFT 阶段为后续 RL 建立了一个高质量的冷启动策略}。

\begin{importantbox}{SFT 的定位}
SFT 不是「比 RL 弱的替代品」，而是 RL 的\textbf{起点}。它负责把通用 agent 能力（工具调用、格式、长任务）装进模型，RL 负责在此基础上把某个领域的能力磨尖。
\end{importantbox}

\subsection{SFT 与 RL 的对照}

课堂把两者的差异压成一张表：

\begin{center}
\begin{tabular}{p{2.4cm}p{5.0cm}p{5.0cm}}
\toprule
 & SFT & RL \\
\midrule
学习来源 & 记录下来的轨迹（不论由谁产生） & 模型自己的 rollout，由奖励打分 \\
信号 & 每个 token 一个目标 & 每条 rollout 一个数 \\
需要 & 数据 + GPU & 活的环境、沙箱、可检查的奖励 \\
擅长 & 学工具调用、格式、长任务 & 把模型已经会尝试的东西磨尖 \\
失效 & 示范本身错误、狭窄或已耗尽 & 起点太弱，或奖励不可检查 \\
\bottomrule
\end{tabular}
\end{center}

\subsection{本章小结}
训练流水线是「预训练 $\rightarrow$ 中期训练 $\rightarrow$ SFT $\rightarrow$ RL」，SFT 是 RL 的冷启动：先装通用 agent 能力，再由 RL 磨领域能力。DeepSeek R1.0 的教训说明跳过 SFT 直接 RL 不可读、不可用；SFT 与 RL 的信号、依赖和失效边界完全不同，二者是接力关系而非替代关系。

\section{轨迹作为 token 序列：模板、mask 与损失}

\subsection{模板把整条轨迹渲染成一个字符串}

语言模型训练只认 token 序列，因此要先把一条对话轨迹用 \textbf{chat template} 渲染成一个字符串：system、user、assistant、observation 之间的边界由模板的特殊 token 标出（例如 Qwen 系列的 \texttt{<|im\_start|>}、\texttt{<|im\_end|>}、tool call 标记等）。一条轨迹里可能有十几次工具调用、最后提交给 verifier 的最终答案，但训练基础设施看到的只是\textbf{一整个字符串}。

\subsection{只让 assistant token 承担损失}

与普通语言模型 SFT 不同，agentic SFT 通常\textbf{不训练整条轨迹}：系统提示、用户消息、工具返回的 observation 是「上下文」，不算损失；只有 assistant 自己写出来的部分才算预测目标。具体做法是给每个 token 配一个 0/1 的 mask，损失写成

$$
\mathcal{L}(\theta) = -\sum_{t:\, m_t=1} \log p_\theta(x_t \mid x_{<t}),
$$

其中 $x_t$ 是第 $t$ 个 token，$x_{<t}$ 是它之前的全部上下文，$p_\theta$ 是模型给出的条件概率，$m_t\in\{0,1\}$ 是 mask 位：1 表示这个 token 参与损失，0 表示只作为上下文。最小化它，就是在提高「给定前文时复现记录动作」的概率——这正是行为克隆（imitation learning）的目标。

为什么不在 observation 上训练？课堂的解释是：观测由用户或环境给出，不是 assistant 产生的，训练它等于让模型去预测「别人会说什么」。不过课堂也补充：确实有工作训练观测以帮助模型学习世界状态，所以这是「一般配方」，值得按场景做消融。

\subsection{mask 是怎么被构建出来的}

mask 的构建分三步：\textbf{渲染}（renderer 用模板把对话写成整串并记下每个 assistant 段的起止）、\textbf{分词}（字符位置变成 token 位置）、\textbf{掩码}（assistant 段内为 1，其余为 0）。这是「模板做记账」的直接后果：训练开始时，system/user/observation 这些角色信息已经消失了，训练器只看到一个字符串和一串 bit。

\subsection{主流框架的 mask 开关与 end-of-turn}

课堂点名了几个 SFT 框架的掩码控制：TRL、Axolotl、LLaMA-Factory 都提供「是否训练 prompt」「是否 mask 历史」「是否训练整条轨迹」之类的开关。实操中最常见的错误是\textbf{分不清哪些该算损失、哪些只是上下文}：例如 assistant 消息的段头（row header）不是模型生成的，不该训练；模型真正生成的是内容段；最后还有一个停止 token 要训练；而 \textbf{end-of-turn 往往由多个事件组合}——结束消息、结束工具调用、结束任务。Kimi K3 的模板把它们拆成独立 token，让模型更容易学会「在哪里停」。

\subsection{reasoning block：训还是不训}

另一个重要开关是 \textbf{reasoning（思维链）块}。reasoning 轨迹通常远长于真正的工具调用，若全部计入损失，损失会被长推理主导，模型可能学会「只推理、不动手」。常见做法有三种：把 reasoning 留在上下文里但不训练（它解释后面的动作）；训练它（模型因此学会自己写这种风格）；或者干脆把它移出上下文。课堂举了 Nemotron 3 Ultra 的例子：它把被预算截断的 reasoning 留在上下文里，但把人为截断的那一段 mask 掉、不计损失。另有工作（Balanced SFT）专门\textbf{平衡 reasoning 与工具调用}两类轨迹，避免一类压过另一类。

\begin{warningbox}{三个常见的 mask 错误}
\begin{enumerate}
\item 把 assistant 段头或模板特殊 token 当成可训练 token；
\item 忘记训练 end-of-turn / stop token，导致模型不会停、把对话一直写下去；
\item 让超长的 reasoning 主导损失，训练出一个「只会想、不会做」的 agent。
\end{enumerate}
\end{warningbox}

\subsection{样本权重不等于 token 权重}

课堂用 MAI（Microsoft AI）的 Consolidation SFT 配方说明一个容易忽略的偏差：他们的数据混合\textbf{按样本数}配比为 STEM 与 Coding 占 56\%、Agentic Capability 占 11\%、General Helpfulness 与 Safety 占 33\%；但损失是\textbf{按 token} 计算的，由于 STEM/Coding 轨迹特别长，最终 token 占比变成 89\% / 9\% / 2\%。

\begin{knowledgebox}{配比检查清单}
如果你按「轨迹条数」平衡数据，还要再检查一次「token 分布」是不是你想要的；反过来，按 token 平衡也可能让短而重要的能力被稀释。两种平衡都只是手段，最终要看评测。
\end{knowledgebox}

\subsection{这一块的开放问题}

课堂明确指出两点：其一，\textbf{没有公开实验}系统改变 agent SFT 的 mask 并测量其影响；其二，关于「是否训练 prompt」的谨慎实验大多发生在\textbf{单轮指令微调}上，而 agent 轨迹恰好相反（长 prompt、短答案、数据少）。因此「该不该在 observation 或 prompt 上训练」在 agent 场景仍是开放问题。

\subsection{本章小结}
一条轨迹先被模板渲染成一个 token 字符串，再用 0/1 mask 只让 assistant token 进损失，这就是行为克隆。mask 由「渲染 $\rightarrow$ 分词 $\rightarrow$ 掩码」三步构建，框架给了各种开关；end-of-turn 与 reasoning 块的处理会显著改变学到的行为；配比要同时看样本与 token。这些细节没有统一最优解，需要按场景做消融。

\section{选轨迹：从哪里来、怎么过滤}

\subsection{轨迹来源的四象限}

课堂把 agent SFT 数据的来源总结成 $2\times 2$：谁跑的任务（更强的模型 / 模型自己）$\times$ 保留全部还是只保留通过验证的轨迹。只保留强模型生成的通过轨迹是\textbf{最常见}的一格（SWE-Gym、Nemotron、OT-Agent 等都在此列）；全部保留叫 plain distillation；模型自己生成且全部保留，即 MAI 的自蒸馏；模型自己生成但只保留通过验证的，就是 \textbf{expert iteration}（后续课程会讲）。

课堂以 SWE-Gym 的结果说明强模型蒸馏的收益：用 GPT-4o、Claude 等强模型跑轨迹、只保留通过仓库测试的轨迹，学生模型的解决率显著提高。

\subsection{成功轨迹里也可能有坏行为}

一个重要提醒：即使最终任务成功，中间也可能包含不该学的东西。课堂展示了两条「通过了单元测试」的轨迹：一条大部分回合都在「改一下、重跑测试」的循环里空转；另一条的最终补丁里留着调试用的 print 语句。SFT 的损失会\textbf{照抄每一个动作}，好坏一起学——于是「会修」和「会乱来」被同时刻进权重。

\subsection{按形状过滤：最小轮数}

最朴素的过滤是\textbf{按轨迹长度/轮数}。课堂引用了 openhands agent 论文的消融：在他们的设置里，最好的单一过滤是\textbf{把少于 5 轮的轨迹全部丢掉}。理由很直接：轨迹太短，要么是一把就做对的简单题（学不到错误恢复），要么是生成过程本身就出了问题。

课堂也诚实指出这种过滤的盲点：轮数只衡量「花了多少功夫」，不衡量「功夫好不好」。一个 4 轮干净修复会被丢掉，一个 14 轮乱七八糟的轨迹反而会被保留——但平均而言它仍然有帮助，这就是工程上的「够用的启发式」。

\subsection{选老师：最强不等于最好教}

课堂用一张教师模型消融表说明：把同一个学生模型分别蒸馏自不同教师，教师的「自身强弱」与「教得好不好」并不一致。表中最强的 GPT-5.3-Codex 反而是最差的老师（平均 11.94、归一化 $-1.47$），比量化后的 GLM 4.7（18.16、$+0.73$）在 Terminal-Bench 2.0 上还低约 5 个百分点；Kimi K2.5 与 GLM 5 居中。

可能的原因：强模型的「思考方式」与学生差异太大，轨迹对学生不可学；排行榜只说明模型自己会做题。工程建议是\textbf{多试几个老师}，并用代价低的小模型快速验证学生收益。

\subsection{任务来源与数据多样性}

课堂强调：不同任务来源对不同 benchmark 的贡献差异很大。例如 SWE-smith 对 SWE-bench Verified 贡献最大，但对 Terminal-Bench 贡献较小——\textbf{任务源的选择可能比其它任何超参都更影响结果}。经验结论是：混合 \textbf{4--8 个来源}最好，铺得更宽反而不再有帮助。

另一个反直觉发现是「数据越多越好」并不总成立：如果只是从同一分布\textbf{上采样}（同一任务的更多 rollout），性能会饱和甚至下降；而增加\textbf{分布多样性}（新任务、合成增广）时，性能还能继续上升。课堂用三张曲线对比了「同分布上采样」与「合成增广」。

\subsection{课堂讨论：三条轨迹怎么选}

课堂给了一个讨论题：同一个任务有三条轨迹——A 是 4 轮就干净修好的短轨迹；B 是长轨迹，中间做错了一次编辑、测试失败后又恢复；C 是试了 11 次、每次都没改对，最后却解决了。该保留哪几条、哪些动作该计损失、最小轮数过滤在这题上会错在哪？

课堂的结论是 \textbf{A 和 B 都要}：B 教的是错误恢复能力（这正是 agent 最稀缺的能力之一），A 教的是干净利落；但 B 也有风险——模型可能只学会「先做错再改」，所以要在数据里保留两端、并控制比例。而「最小 5 轮」的过滤在这题上会\textbf{丢掉 A、留下 C}，恰好展示了启发式过滤的代价。

\subsection{本章小结}
轨迹来源要看清「谁生成、是否过滤」；成功轨迹也可能带坏行为，需要按形状/规则过滤；教师不是越强越好，要实测学生收益；任务源与数据多样性比单纯堆数据重要（4--8 个来源、增广优于上采样）；课堂讨论说明「保留正确 + 恢复两类轨迹」比机械按轮数过滤更接近想要的 agent 能力。

\section{统一格式与多样性：Agent Data Protocol}

\subsection{异构数据集带来的工程墙}

课堂介绍了主讲人自己的一个工作（与 Graham Neubig 教授合作）：\textbf{Agent Data Protocol}。动机很实际——现存的 agentic SFT 数据集由不同团队、针对不同目标、用不同格式构建：同样做 web 任务，有的用 HTML 作为观测，有的用 accessibility tree；工具调用的表示也各不相同。如果想把 GoBrowse 与 Mind2Web 一起训练，就得为每个数据集写一份到目标 agent 格式的转换器；换成 OpenHands，又要重写一遍。数据集一多，转换器数量是「数据集 $\times$ 目标 agent」的乘积。

\subsection{中间格式：每个数据集只转一次}

Agent Data Protocol 的思路是引入一个\textbf{统一的标准格式}：每个数据集只需要转换一次到标准格式，再从标准格式转换一次到目标 agent 格式。web 数据集中，HTML 与 accessibility tree 两种观测都保留在统一结构里，下游可以自选；这样既减少了工程量，又让多数据集混训成为可能。转换之后仍需检查：tool call 是否都能解析、每个调用是否带 function schema、对话是否按预期结束。

\subsection{多源混合确实比单源好}

论文做了对照实验：固定同一个 harness、同一个模型、同一套评测，只改变训练数据的混合。结论是：\textbf{几乎每个 agent benchmark 上，多源混合都优于单源训练}，甚至在该领域自己的 benchmark 上，「混合」也能打败「恰好匹配该领域的单源数据」。这说明数据多样性带来的收益大于「领域精确匹配」的收益。

\subsection{本章小结}
Agent 数据格式高度异构，逐对写转换器不可扩展；统一中间格式把「数据集 $\times$ agent」的乘积降为「数据集 + agent」两段线性转换。多源混训在几乎所有 benchmark 上优于单源，多样性是 agent SFT 的重要杠杆。

\section{跑通一次 SFT：工程细节}

\subsection{来自真实技术报告的配置}

课堂列了几份真实配置作参照：Nemotron 3 Ultra 把序列打包到约 294K token、batch size 64，学习率随时间衰减；MAI 的自蒸馏用 128K 序列、很大的 batch，并在正式 SFT 前做 warm-up；另一个近期开源模型用 512K 序列、分几个阶段，学习率在高量子集上衰减。课堂的建议是：动手前至少读一两份前沿实验室的技术报告，里面有很多可以借鉴的细节。

\subsection{MoE 的专家路由}

前沿模型大多是 MoE（混合专家）。课堂指出一个现象：在狭窄的 agent 数据上微调 MoE 时，\textbf{少数专家会吃掉大部分 token}（图右侧：专家序号越大，分到的 token 越少）。MAI 在 SFT 阶段对路由做很强的平衡，难度比 RL 阶段高约一千倍，并把 dropout 提到 0.15。

\begin{knowledgebox}{为什么 SFT 阶段的路由更容易失衡}
SFT 数据窄、格式单一，路由器会迅速把同类 token 都送给同一个专家，形成「赢者通吃」。一旦路由塌缩，模型容量利用率下降，后续 RL 也会被限制。因此 MoE 的 SFT 往往需要辅助的负载均衡损失或路由约束。
\end{knowledgebox}

\subsection{打包整段对话，不截断}

打包（packing）的常见做法是把多条对话塞进固定长度的序列里以提高利用率，但 agent SFT 有一个额外约束：\textbf{不要截断、不要切分轨迹}，要保留完整对话。课堂给出的做法是：每条对话被分配到「剩余容量最贴合它」的那个 pack，任何对话都不被切开；相同问题的对话不放进同一个 pack，pack 之间再随机打乱。这样做的一个动机是\textbf{避免幻觉}——被切断的轨迹会让模型学到错误的续写边界。

\subsection{模板漂移（template drift）}

一个容易被忽视却致命的工程细节是：\textbf{训练用的模板必须与部署用的模板完全一致}，甚至空白字符都要一致。课堂给了一个具体场景：同一条轨迹，在「工具消息已经追加」前后渲染两次，会得到两个不同的字符串；一旦条件性的 thinking block 改变了前面 assistant 段的结尾，前缀就不再匹配。在\textbf{工具循环}里前缀不匹配是致命的——模型在训练时见到的上下文与推理时见到的不是同一个东西，行为会不可预期地退化。

\begin{warningbox}{模板漂移是「静默失败」}
它不会让 loss 报错，只会让训练出来的能力在部署时打折。上线前务必用同一个 renderer 分别渲染「训练样本」和「真实推理请求」，逐 token 对比。
\end{warningbox}

\subsection{公开的训练工件}

课堂推荐关注公开发布训练细节的模型：有的 model card 给出分阶段训练表（每阶段的 token 数、序列长度、目的），公开 W\&B 上的训练损失曲线，甚至放出中间 checkpoint 供继续训练或后训练；也有模型同时发布后训练 checkpoint、训练数据与配方。对想学习「工业级 SFT 长什么样」的人来说，这些工件比论文正文更有信息量。

\subsection{本章小结}
真实 SFT 的工程点集中在：序列打包不截断、MoE 路由平衡、模板与部署严格一致、学习率随阶段衰减。技术报告与公开工件是学这些细节的最佳材料；其中模板漂移是最容易静默毁掉一次训练的坑。

\section{训练好不等于评测好：训练与部署的失配}

\subsection{训练只走记录过的动作}

训练与推理有一个结构性差异：训练时，模型在\textbf{记录的轨迹}上一步一步算损失，每一步的上下文都是真实发生过的；而部署时\textbf{没有记录}，模型每做一个动作就会改变接下来能看到的状态。课堂的例子是：模型做了一次与训练数据不同的编辑，测试失败，于是进入一个「训练数据从未访问过的状态」；此时怎么做，取决于它的错误恢复能力，而这正是 SFT 数据里可能没有的。

因此课堂给了一个重要判断：\textbf{训练损失曲线可以用来监控训练是否正常，但不能替代真实评测}。低损失只保证「在记录过的状态上能复现动作」，不保证「在没见过的状态上做对」。

\subsection{评测的可信度清单}

课堂给了一串很实用的评测建议：

\begin{enumerate}
\item \textbf{先过拟合 20 个样本}：如果 loss 降不下去，bug 大概率在上游（数据、模板或 mask）；
\item \textbf{用 harness 跑有可检查结果的任务}，而不是只看训练 loss；
\item \textbf{小心读论文数字}：有些数字是「多个 harness 里最好的那个」，不等于你单次单 harness 能复现的结果；
\item \textbf{让 split 匹配你要主张的结论}：要声称「泛化到新仓库」，就要按 repository 切分，并保证同一任务的所有 run 在同一侧；
\item \textbf{跨 harness 收集轨迹}：在一个系统里收集的轨迹会教模型「这个系统的约定」，换到 Claude Code 或 Codex 可能不泛化；
\item \textbf{平衡按 token，也要评测你正在改进的能力}，不要只在一个领域上 bench-max。
\end{enumerate}

\subsection{本章小结}
训练沿记录轨迹、部署要自己走路，二者失配意味着 loss 不能替代评测。评测要先用过拟合小样本做 sanity check，再用 harness 跑可检查任务；读数字要对齐 split 与 harness，并跨 harness 收集数据以防「只会在一个系统里工作」。

\section{SFT 与 RL 的交接：SFT 的收益在哪里}

\subsection{三条经验}

课堂给出三条经验（以 DeepSeek R1.0 与 OpenThoughts 系工作为据）：第一，\textbf{完全不做 SFT、直接对 base model 做 RL，几乎带不动 base model}；第二，\textbf{只做 SFT 也远达不到 SFT+RL 的高度}；第三，最佳配方是从一个「\textbf{故意没有训到自身最高分、而是为后续 RL 训得更好}」的 checkpoint 开始 RL。

课堂用一张表展示「SFT 与 RL 的配比」：ColdSFT+RL 的 8B 模型在 SWE-bench Verified（100）、OT-TBLite、Terminal-Bench 2.0 上分别达到 35.7、16.0、13.5，平均 21.7；只做 SFT 的版本平均 15.1--15.9；完全不 SFT 直接 RL 只有 3.6；base 模型 2.7。课堂的总结句是：\textbf{「RL 的收益最大，当 SFT 模型是带着 RL 的目标选出来的。」}

\subsection{没有更强的教师怎么办}

课堂最后的问答讨论了一个很实际的问题：如果没有比学生更强的模型来生成轨迹，怎么办？课堂给了两条路：一是像 MAI 那样\textbf{自蒸馏}（用模型自己的推理轨迹）；二是像一些前沿实验室那样，把 SFT 压缩成「冷启动」，主要依靠 \textbf{RL} 提升能力。课堂的判断是：SFT 的天花板受教师能力限制，如果没有更好的教师，就必须更多依赖 RL。

\subsection{要不要监督错误答案}

另一个问答问到「是否要在错误答案上做监督」。课堂的回答是分两层的：\textbf{中间的错误动作值得监督}（模型要学「做错了怎么恢复」），但\textbf{最终结果必须正确}；也就是说，监督的是「错误 + 恢复 + 正确结果」的完整过程，而不是「做错并停在错处」。这和前面「成功轨迹里也可能有坏行为」的提醒互相呼应：数据筛选要保留恢复能力、剔除纯粹的坏习惯。

\subsection{本章小结}
SFT 提供通用 agent 能力的冷启动，RL 在它之上磨领域能力；「不 SFT 直接 RL」几乎无效，而「SFT 训到最高分再 RL」也不是最优——要给 RL 留空间。没有更强教师时靠自蒸馏或 RL；数据上要监督「错误 + 恢复」，但最终结果必须正确。

\section{总结与延伸}

本讲把 SFT 放回完整训练流水线中：它是 RL 的冷启动，负责把通用 agent 能力（工具调用、格式、长任务）写进权重。全讲可以压缩为六条：

\begin{enumerate}
\item 更新 agent 有三条路（上下文、外部工件、权重）；改权重换来快推理与彻底行为改变，代价是训练成本、黑盒与模型绑定。
\item 轨迹先被 chat template 渲染成一个 token 串，再用 0/1 mask 只让 assistant token 进交叉熵——SFT 本质是行为克隆。
\item mask 的细节决定学到什么：end-of-turn、reasoning 块、样本权重 vs token 权重，都需要显式控制；这些没有统一答案，要做消融。
\item 数据决定上限：来源（强模型/自己、是否过滤）、过滤（最小轮数等）、教师选择（最强 $\neq$ 最好教）、任务来源（4--8 个混合、增广优于上采样）。
\item 工程细节会静默毁掉训练：打包不截断、模板不许漂移、MoE 路由要平衡；公开技术报告与训练工件是最好的教材。
\item 训练 loss 不能替代评测：训练沿记录轨迹、部署要自己走路；先用 20 个样本过拟合做 sanity check，再对齐 harness 与 split。
\end{enumerate}

\begin{importantbox}{一句话记住本讲}
\textbf{SFT 把「记录下来的轨迹」变成「权重里的行为倾向」；它决定 agent 的起点，但不决定终点——终点要靠 RL，而起点要专门为 RL 挑。}
\end{importantbox}

\textbf{延伸思考}：(1) 把 mask 当作一种「数据加权」来看，agent SFT 里所有关于 reasoning/tool call/observation 的取舍，都是在决定「让哪些 token 塑造模型」；(2) 教师选择与数据来源选择的收益，可能比调学习率大一个量级，这意味着 agent SFT 的核心工作更像是「数据工程 + 评测设计」，而不是「训练脚本调参」；(3) 模板漂移、打包截断这类问题与「评测口径」问题同源——它们都在问「训练时看到的世界，和部署时的世界是不是同一个」。下一讲进入强化学习基础，讨论如何用奖励而不是示范来继续提升 agent。

\end{document}
